% Short answers to Chapter 8's exercises, lectures/L08/appendix/b_exercises.md. Exercises 8.3 and
% 8.4 are VHDL only and have no entry.

\answerchapter{c8:ch}

\answer{c8:ex:1}{A binary-count Moore machine}
\pt{a}Four bubbles, \code{00/0}, \code{01/0}, \code{10/1} and \code{11/1} as \code{Q1Q2/Y}, in a ring
\code{00 -> 01 -> 10 -> 11 -> 00} with one arrow per state labelled \code{X = 1}, a self-loop labelled
\code{X = 0} on each, and reset entering \code{00}.
\pt{b}\begin{tabular}[t]{@{}l@{\quad}*{8}{c}@{}}
  \code{Q1Q2X} & 000 & 001 & 010 & 011 & 100 & 101 & 110 & 111 \\
  \midrule
  \code{Q1+} & 0 & 0 & 0 & 1 & 1 & 1 & 1 & 0 \\
  \code{Q2+} & 0 & 1 & 1 & 0 & 0 & 1 & 1 & 0 \\
  \code{Y} & 0 & 0 & 0 & 0 & 1 & 1 & 1 & 1 \\
\end{tabular}
\pt{c}\code{Q1+ = Q1X' + Q1Q2' + Q1'Q2X}: two pairs and the lone cell \code{011}, all three
essential (it is \code{Q1} XOR \code{Q2X}). \code{Q2+ = Q2X' + Q2'X}, two pairs, which is
\code{Q2 XOR X}. \code{Y = Q1}, one group of four.
\ptnote{Tip}With the LED lit only in \code{10}, \code{Y = Q1Q2'} is decoded from both bits, and on the
\code{11 -> 00} transition, if \code{Q2} falls before \code{Q1}, the state reads \code{10} for an
instant and \code{Y} glitches high. If \code{Y} drove something that cared, register it: a flip-flop
loaded with \code{Q1+Q2+'}, so \code{Y} changes only on the clock edge.

\answer{c8:ex:2}{A Moore \code{1, 0, 1} detector}
\pt{a}Four states: \code{S0} nothing matched, \code{S1} saw \code{1}, \code{S2} saw \code{10}, and
\code{S3} saw \code{101}, the only state with \code{y = 1}; reset enters \code{S0}. Its arrows are the
rows of the table in \textbf{b)}.
\pt{b}\begin{tabular}[t]{@{}l@{\quad}*{8}{c}@{}}
  state, \code{din} & \code{S0},0 & \code{S0},1 & \code{S1},0 & \code{S1},1 & \code{S2},0 &
    \code{S2},1 & \code{S3},0 & \code{S3},1 \\
  \midrule
  next state & \code{S0} & \code{S1} & \code{S2} & \code{S1} & \code{S0} & \code{S3} & \code{S0} &
    \code{S1} \\
  \code{y} & 0 & 0 & 0 & 0 & 0 & 0 & 1 & 1 \\
\end{tabular}
\pt{c}For example, \code{din = 1, 1, 0, 1, 0, 1}:\par\smallskip
\noindent\begin{tabular}{@{}l@{\quad}*{6}{c}@{}}
  cycle & 1 & 2 & 3 & 4 & 5 & 6 \\
  \midrule
  state before & \code{S0} & \code{S1} & \code{S1} & \code{S2} & \code{S3} & \code{S0} \\
  \code{din} & 1 & 1 & 0 & 1 & 0 & 1 \\
  \code{y} & 0 & 0 & 0 & 0 & 1 & 0 \\
\end{tabular}\par\smallskip
\noindent\code{y} is high in cycle 5, the cycle after the match's last bit; the \code{1, 0, 1} in
bits 4 to 6 overlaps the match and is correctly not reported.
\pt{d}A \code{0} in \code{S0}, \code{S2} or \code{S3} returns to the start. A \code{1} in
\code{S1} must stay in \code{S1}, not fall back to \code{S0}, because that \code{1} is itself the
first bit of a possible match (\code{1101} ends in \code{101}); likewise a \code{1} in \code{S3} goes
to \code{S1}, since non-overlapping discards the match's bits but not a new one. A \code{0} in
\code{S2} goes all the way back, because \code{100} ends in no prefix of \code{101}.

\answer{c8:ex:5}{Accounting for \code{fsm_led}}
\pt{a}A timeout comes every \code{TICK_COUNT + 1} = 5,000,001 cycles of 20\,ns, 100\,ms. The blink
frequency is 5\,Hz: the LED toggles once per timeout, so one on-and-off cycle takes two timeouts,
200\,ms, and the 10\,Hz you might first write down is the toggle rate.
\pt{b}The override is \code{constant TIMER_TICK_COUNT: natural := 3;}, passed in by
\code{generic map(TIMER_TICK_COUNT)}. One real blink period is 2 $\times$ 5,000,001, about ten million
clock cycles: 200\,ms of simulated time at 50\,MHz, far past the 10\,ms \code{--stop-time}, and
tens of thousands of times longer than the whole testbench takes with the override.
\pt{c}\begin{itemize}
  \item From \code{STATE_OFF}, three presses of \code{button_n(0)} visit \code{STATE_BLINK},
    \code{STATE_ON} and \code{STATE_OFF} in turn; three presses of \code{button_n(1)} visit
    \code{STATE_ON}, \code{STATE_BLINK} and \code{STATE_OFF}.
  \item The machine stays where it is. The two lines are the assignments to \code{to_prev_state} and
    \code{to_next_state}, each of which requires the other button's edge to be \code{'0'}.
  \item The timer resumes from wherever it was frozen, so the first toggle after re-entry comes
    anywhere from one cycle to a full period later, and the first blink interval is short. Holding
    \code{timer_enable} high is neither better nor worse for correctness, since re-entry would still
    catch the timer at an arbitrary point in its period; it only keeps a counter switching, and
    spending power, when nothing reads it.
\end{itemize}

\answer{c8:ex:6}{The serial receiver}
\pt{a}\begin{itemize}
  \item \code{STATE_IDLE}: leaves on a tick with \code{rx_s2 = '0'} and the previous tick's
    \code{rx_s2 = '1'}, clearing \code{ticks}; the history is one flip-flop loaded with
    \code{rx_s2} on every tick.
  \item \code{STATE_START}: waits 8 ticks; then \code{rx_s2} low goes to \code{STATE_DATA},
    clearing \code{ticks}, and high goes back to \code{STATE_IDLE}.
  \item \code{STATE_DATA}: waits 16 ticks per bit, eight times; each sample shifts \code{rx_s2} into
    \code{frame} and clears \code{ticks}, and the eighth moves on to \code{STATE_STOP}.
  \item \code{STATE_STOP}: waits 16 ticks; then goes to \code{STATE_IDLE}, pulsing
    \code{byte_valid} with \code{data_out <= frame} if \code{rx_s2} is high, and \code{frame_err} if
    it is low.
  \item Every state has a self-loop for ``no tick yet, or not the right count yet''.
\end{itemize}
\pt{b}\begin{itemize}
  \item 50,000,000 / 9600 = 5208.33 cycles per bit, and 325.52 per sixteenth.
  \item \code{TICK_COUNT} = 324.52. 324 gives 325-cycle ticks, 0.52 cycles short ($-0.16\%$); 325
    gives 326, 0.48 cycles long ($+0.15\%$). The default is 325, the closer of the two.
  \item 9.5 bit periods is 152 ticks: 152 $\times$ 326 = 49,552 cycles against 49,479, about 73 cycles
    late, 0.22 of a tick. The margin was 8 ticks, half a bit.
  \item The drift within one frame would have to reach the margin: at 0.024 of a tick per bit,
    that is a frame of some 340 bit periods (less once the crystals' own tolerance adds to it).
\end{itemize}
\pt{e}\begin{itemize}
  \item The low-stop-bit frame is handed over as a byte and \code{frame_err} never pulses. The
    first check to fire is ``a frame with a low stop bit must raise \code{frame_err}!'', which the
    testbench makes before its byte-count check.
  \item Add a case that pulls \code{rx} low on an idle line for a few ticks, well under half a bit,
    then holds it high for longer than a frame, and asserts that neither \code{valid_count} nor
    \code{error_count} changes. Without the re-check the glitch starts a frame of all ones with a
    high stop bit, and \code{x"FF"} is handed over as a byte.
  \item The testbench's transmitter runs off the receiver's own clock at exactly sixteen ticks a
    bit, so nothing drifts, and even its noisy frame holds each bit clean from tick 3 to tick 13.
    What is given away is the margin on the early side, down from 8 ticks to 3 or less, so a
    slightly slow transmitter or a slow edge reads the previous bit.
  \item A simulator models every signal as a clean \code{'0'} or \code{'1'} changing at an exact
    instant, and every flip-flop as sampling perfectly; it has no setup and hold window, so it
    cannot produce metastability (and the testbench changes \code{rx} just after a clock edge
    anyway). The failure is physical and probabilistic, so only the rule prevents it.
  \item It promises that for the stimuli the testbench applies, in an ideal simulation, the outputs
    are the ones it checks for. It does not promise anything about inputs it never applies, such as
    glitches or drifting clocks, about timing margin, or about what the simulator does not model:
    metastability, glitches and timing on the real chip.
\end{itemize}

\answer{c8:ex:7}{The serial transmitter}
\pt{a}Idle: \code{tx = '1'}; on \code{send = '1'}, latch \code{data_in} into \code{frame} (the
byte is captured here) and start the timer. Start bit: \code{tx = '0'} for one timeout. Data:
\code{tx = frame(7)}, shifting \code{frame} up one place per timeout, eight times. Stop:
\code{tx = '1'} for one timeout, then idle. \code{busy} is \code{'1'} whenever the state is not
idle.
\pt{e}\begin{itemize}
  \item \code{busy} drops at the end of the eighth data bit, so a producer that sends the next byte
    straight away puts its start bit right after the last data bit, with no stop bit. The receiver
    samples its stop bit inside that start bit and rejects the frame with \code{frame_err}, and if
    the last data bit was \code{0} the new start bit has no falling edge for it to find. It first
    shows up on the first pair of bytes sent back to back; one byte followed by idle looks perfect.
  \item The check ``\code{busy} must stay asserted for the whole stop bit!'', made in the middle of
    the first frame's stop bit. On \code{tx} a skipped stop bit is idle, which is the same high
    level, so the line is identical until another frame follows; only \code{busy} shows that the
    transmitter thinks the frame is over.
  \item The sample point moves 1\% of a bit per bit and starts half a bit from either edge, so it
    reaches the wrong bit about 50 bits in (nearer 44 counting the tick of start-edge uncertainty).
    A ten-bit frame drifts only a tenth of a bit, so a start bit per byte resynchronizes long before
    the drift matters, where a single start bit for a long message would lose it within a few
    bytes.
\end{itemize}
